\scanpage{12}
% archive foliation: 11
(où $X\xrightarrow[r]{u}Y\xrightarrow[s]{v}Z$ sont des morphismes $\in\mathcal N$, de degré $r,s$) :
\diagram{\begin{tikzcd}Y\arrow[r]&\Sigma^r(X/Y)\end{tikzcd}}
\diagram{\begin{tikzcd}[column sep=large,row sep=large]
Z\arrow[r,"{\unsure{\operatorname{Cor}|\operatorname{Cn}}_{Y/Z}}"']\arrow[rrd,"{\deleted{\unsure{\longrightarrow}}}" description]\arrow[rrdd,"{\unsure{\operatorname{Cor}|\operatorname{Cn}}_{X/Z}}"]\arrow[rrddd,"{\deleted{\unsure{\longrightarrow}}}" description]&\Sigma^s(Y/Z)\arrow[r,"{\Sigma^s(\unsure{\operatorname{Cor}|\operatorname{Cn}}_{X/Y}/Z)}"]&\Sigma^s(\Sigma^r(X/\unsure{Y})/Z)\arrow[d,"{\mathrm{can}}"]\\
&&\Sigma^s(\Sigma^r(\deleted{\unsure{X}}X/Z)/Z)\arrow[d,"{\mathrm{can}}"]\\
&&\Sigma^{sr}(X/Z)\\
&&{}
\end{tikzcd}}

(iii ter) Soit $u:X\longrightarrow Y$ dans $\mathcal N$, et $Y'\longrightarrow Y$ quelconque, alors $\unsure{\operatorname{Cor}|\operatorname{Cn}}_{X'/Y'}\deleted{Y'\longrightarrow}$ est ``image réciproque'' de $\unsure{\operatorname{Cor}|\operatorname{Cn}}_{X/Y}$ par $Y'\longrightarrow Y$, i.e. le diagramme
\diagram{\begin{tikzcd}
\Sigma^r(X/Y)&Y\arrow[l,"{\unsure{\operatorname{Cor}|\operatorname{Cn}}_{X/Y}}"']\\
\Sigma^r(X'/Y')\arrow[u,"{\mathrm{can}}"]&Y'\arrow[l,"{\unsure{\operatorname{Cor}|\operatorname{Cn}}_{X'/Y'}}"']\arrow[u]
\end{tikzcd}}
est commutatif

(iv ter) Soient $u_i:\deleted{\unsure{X}}X_i\longrightarrow Y\quad(i=1,2)$ dans $\mathcal N$. Soit $X=X_1\amalg X_2$, $u=(u_1,u_2)$ ; donc pour $\forall$ entier $r$, on a
\[
\Sigma^r(X/Y)=\deleted{\unsure{\Pi}}\coprod_{\substack{i+j=r\\i,j\geq0}}\Sigma^i(X_1/Y)\times\Sigma^j(X_2/Y)
\]
en particulier, si $r_i=\deg u_i$ et $r=r_1+r_2=\deg u$, on a un morphisme \inserted{\unsure{d'injection|d'immersion}} canonique $\Sigma^{r_1}(X_1/Y)\times\Sigma^{r_2}(X_2/Y)\longrightarrow\Sigma^r(X/Y)$. Ceci dit, on a commutativité du
\diagram{\begin{tikzcd}[column sep=large]
Y\arrow[r,"{(\unsure{\operatorname{Cor}|\operatorname{Cn}}_{X_1/Y},\unsure{\operatorname{Cor}|\operatorname{Cn}}_{X_2/Y})}"]\arrow[rd,"{\unsure{\operatorname{Cor}|\operatorname{Cn}}_{X/Y}}"']&\Sigma^{r_1}(X_1/Y)\times\Sigma^{r_2}(X_2/Y)\arrow[d]\\
&\Sigma^r(X/Y)
\end{tikzcd}}

\uline{Exemple I} Soit $G$ un groupe, soit $C$ la catégorie des ensembles où $G$ opère. Un morphisme $u:X\longrightarrow Y$ de $C$ sera \deleted{dans $\mathcal N$ s'il existe} \inserted{dit de degré $r$} \deleted{\unsure{$r$ tel que pour $\forall$}} si pour $\forall\ y\in Y$, \deleted{\unsure{il y a}} \deleted{\illegible{2}} $\operatorname{card}(u^{-1}(y))=r$. Un morphisme est dans $\mathcal N$ s'il existe $r$ tel qu'il soit de degré $r$. Les conditions (i) ; (v) et (i bis) ; (v bis) sont trivialement vérifiées. [N.B. les produits fibrés et sommes à ens d'indices quelconque existent dans $C$, \deleted{\unsure{et}} de $=$ le quotient d'un objet de $C$ par un groupe d'automorphismes $\unsure{G|H}$]. Soit $u:X-Y$ de degré $r$. Notons que $\Sigma^r(X/Y)$ est l'ensemble \inserted{comme ens. sous-jacent} \deleted{\unsure{des $r$-uplets}} réunion des $\Sigma^r(u^{-1}(y))\quad(y\in Y)$ ; comme $u^{-1}(y)$ a exactement $r$ \unsure{éléments}, \unsure{Alors}, \deleted{\unsure{il}} il y a dans $\Sigma^r(u^{-1}(y))$ un élément privilégié, \unsure{(image} du $r$-uplet \inserted{mod $\mathfrak{S}_r$} de ces $r$ éléments rangés dans un ordre quelconque). D'où une application-section $Y\longrightarrow\Sigma^r(X/Y)$, qui évidemment permute aux automorphismes de la situation, donc : \unsure{fonctorielle}, \unsure{i.e.} \deleted{\unsure{un}} c'est un \deleted{\unsure{hom}} hom de $C$. Si on
